% !PRINTED 239
\weberpage{239}
Es ist aber
\[
\varphi_2^{-1} \varphi_1 \varphi_2 = \begin{pmatrix} 0, & e^{\frac{\pi i}{4}} \\ e^{-\frac{\pi i}{4}}, & 0 \end{pmatrix}
\]
\[
\varphi_1^{-1} \varphi_2 \varphi_1 = \begin{pmatrix} 0, & e^{\frac{\pi i}{4}} \\ e^{-\frac{\pi i}{4}}, & 0 \end{pmatrix},
\]
\[
\begin{aligned}
&c\varphi_1 = \varphi_1 d,\quad d\varphi_1 = \varphi_1 c,\quad cd\varphi_1 = \varphi_1 dc\\
&c\varphi_2 = \varphi_2 dc,\quad d\varphi_2 = \varphi_2 d,\quad cd\varphi_2 = \varphi_2 c,
\end{aligned}
\]
und aus diesen Formeln folgt, dass die Gruppen (11) durch Transformation mit $\varphi_1$ und $\varphi_2$ in $D_2'$ transformirt werden.

Bei der Bestimmung der übrigen Poly\"edergruppen zweiter Art sind die folgenden allgemeinen Bemerkungen von Nutzen.
Die Inversion $j = \begin{pmatrix} i, & 0 \\ 0, & i \end{pmatrix}$ ist eine Aehnlichkeitssubstitution (\S.\,37) und daher mit jeder anderen Substitution vertauschbar, und folglich können wir aus jeder Poly\"edergruppe erster Art wenigstens eine Gruppe zweiter Art herleiten:
\weq{12}{P + Pj.}
Um die anderen etwa noch vorhandenen Gruppen dieser Art zu finden, erinnern wir uns, dass nach dem \gesperrt{Sylow}'schen Satze (\S.\,29, I.) in jeder Gruppe $G$ eine Gruppe enthalten sein muss, deren Grad die höchste Potenz von 2 ist, die im Grade von $G$ aufgeht. Die erweiterten Tetra\"eder-, Octa\"eder- und Ikosa\"edergruppen haben aber die Grade 24, 48, 120, müssen also einen Theiler vom Grade 8, 16, 8 enthalten.

Es sei nun $T$ die Tetra\"eder-, $O$ die Octa\"eder- und $J$ die Ikosa\"edergruppe. In $T$ und $J$ ist eine Vierergruppe $D_2$ enthalten, aber keine Substitution vierter Ordnung, in $O$ eine Di\"odergruppe $D_4$ und keine Substitution achter Ordnung, und es muss daher unter den erweiterten Poly\"edergruppen eine der unter III. betrachteten erweiterten Di\"edergruppen enthalten sein.
% [sic] printed with \"o here (Di\"odergruppe) vs \"e elsewhere on this page (Di\"edergruppen) — reproduced as read.

Von den erweiterten Di\"edergruppen entsteht aber $D_2''$ aus $D_2$ durch Inversion, während
\[
\begin{aligned}
D_2' = &\begin{pmatrix}1, & 0\\ 0, & 1\end{pmatrix},\ \begin{pmatrix}\sqrt{i}, & 0\\ 0, & +i\sqrt{i}\end{pmatrix},\ \begin{pmatrix}i, & 0\\ 0, & -i\end{pmatrix},\ \begin{pmatrix}i\sqrt{i}, & 0\\ 0, & +\sqrt{i}\end{pmatrix},\\
&\begin{pmatrix}0, & -1\\ 1, & 0\end{pmatrix},\ \begin{pmatrix}0, & -i\sqrt{i}\\ \sqrt{i}, & 0\end{pmatrix},\ \begin{pmatrix}0, & i\\ i, & 0\end{pmatrix},\ \begin{pmatrix}0, & -\sqrt{i}\\ i\sqrt{i}, & 0\end{pmatrix}
\end{aligned}
\]
%CONT>>
